\section*{Esercizi proposti}

\begin{enumerate}
  \item Sia $\mathcal{L}(V)$ lo spazio degli endomorfismi di $V$, si consideri $\mathcal{F} = \{ f \in \mathcal{L}(V) : \text{Ker}(f) \subseteq \text{Im}(f) \}$ e sia $\mathcal{A} = \{ f \in \mathcal{L}(V) : V = \text{Ker}(f) \, \oplus \, \text{Im}(f) \}$.
  \begin{enumerate}
    \item Trovare un $A$ appartenente a $\mathcal{F}$ e $B$ non appartenente. Mostrare che $\mathcal{F}$ è un sottospazio oppure dimostrare che non lo è.
    \item Trovare un $A$ appartenete a $\mathcal{A}$ e $B$ non appartenente.
  \end{enumerate}
  \begin{proof}[Svolgimento]
    Ricordiamo che $\text{Ker}(f) = \{ \underline{v} \in V : f(\underline{v}) = \underline{0} \}$ e $\text{Im}(f) = \{ \underline{w} \in W \big| \exists \underline{v} \in V : f(\underline{v}) = \underline{w} \}$, quindi si ha
    che una possibile funzione è isomorfa a questa matrice qua
    \begin{equation*}
      A = \begin{pmatrix}
          0 & 1 \\
          0 & 0
      \end{pmatrix}
    \end{equation*}
    Vediamo che il kernel di questa matrice qua è rappresentata da vettori che appartengono a $\text{Span}(\underline{e_1})$, quindi preso un vettore $\underline{v} = \alpha \underline{e_1} + \beta \underline{e_2}$, quindi si ha che
    $f(\underline{v}) = f(\alpha \underline{e_1} + \beta \underline{e_2}) = \alpha f(\underline{e_1}) + \beta f(\underline{e_2}) = \beta f(\underline{e_2}) = \beta \underline{e_1}$, dunque abbiamo che $\text{Ker}(f)  \subseteq \text{Im}(f)$.
    La matrice $B$ che non appartiene invece a $\mathcal{F}$ è invece la matrice nulla, per esempio, siccome si ha che $\text{Ker}(f) = V$ e, chiaramente, $\text{Im}(f) = \{ \underline{0} \}$, dunque non abbiamo che $V = \text{Ker}(f) \not\subseteq \text{Im}(f)$. Segue necessariamente che $\mathcal{F}$ non è un sottospazio. \\
    Un esempio di matrice $A$ che appartiene a $\mathcal{A}$ è invece pari a
    $$
      A = \begin{pmatrix}
        1 & 0 \\
        1 & 0
      \end{pmatrix}
    $$
    siccome abbiamo che $\underline{e_2}$ in $\text{Ker}(A)$ e preso un vettore $\underline{v} \in \text{Im}(f)$ abbiamo che esso sarà della forma $\underline{v} = \alpha \underline{e_1} + \beta \underline{e_2}$.
    Per quanto riguarda $B$ abbiamo che la matrice
    $$
      B = \begin{pmatrix}
        1 & -1 \\
        1 & -1
      \end{pmatrix}
    $$
    non appartiene a $\mathcal{A}$ siccome si ha che $\text{ker}(B) = \text{Span}(\begin{pmatrix} 1 \\ 1 \end{pmatrix})$ mentre, presa la base $\mathcal{B} = \{ \begin{pmatrix} 1 \\ -1 \end{pmatrix}, \begin{pmatrix} 1 \\ 1 \end{pmatrix} \}$ abbiamo che
    $$
      B(\underline{v}) = B(\alpha \begin{pmatrix} 1 \\ -1 \end{pmatrix} + \beta \begin{pmatrix} 1 \\ 1 \end{pmatrix}) = \alpha B\begin{pmatrix} 1 \\ -1 \end{pmatrix} + \underline{0} = \alpha \begin{pmatrix} 1 \\ 1 \end{pmatrix} \in \text{Ker}(B). 
    $$
    Si ha quindi che $\text{Im}(f) = \text{Ker}(B)$, dunque abbiamo che non possono essere in somma diretta, dunque $B \not\in \mathcal{A}$. Questo insieme non è un sottospazio: infatti, se prendiamo due applicazioni "invertite" (ovvero prese $A, B$ se $B$ mi manda tutti i vettori nel kernel di $A$ in loro stessi allora $A+B$ mi rompe la somma diretta), siccome
    $$
      A + \begin{pmatrix} 0 & -1 \\ 0 & -1 \end{pmatrix} = B,
    $$
    quindi questo insieme non è chiuso rispetto alla somma.
  \end{proof}
  \item Trovare la matrice inversa
\end{enumerate}
